\chapter{Centered Ramp Sums and the Window Landau Theorem}
\label{ch:02}

The classical theorem of Landau~\cite{Landau1905} asserts that a Dirichlet series with nonnegative coefficients has a singularity at the real point on its abscissa of convergence. For the logarithmic derivative $-L'(s, \chi)/L(s, \chi)$, nonnegativity fails for complex characters $\chi$ and holds only when combined into a positive linear combination with $-\zeta'(s)/\zeta(s)$; moreover, the automorphic period tower of Parts~II--VI does not supply an asymptotic bound on all sufficiently large $y$, nor even a two-sided bound—it supplies only a \emph{one-sided upper bound} $G_{\mathcal{F}}(y) \le C y^\theta$ across a sequence of \emph{arbitrarily large exponential windows} $[Y, Y^\lambda]$ ($\lambda > 1$, specifically $\lambda = 3/2$).

In this chapter we prove from first principles the \emph{Window Landau Theorem} (\cref{thm:02:window-landau}): a one-sided upper bound on a positively weighted admissible family of centered prime-power ramp sums across arbitrarily large exponential windows $[Y, Y^\lambda]$ forces all $L$-functions in the family to be nonvanishing throughout the half-plane $\Re(s) > \theta$. Every step is completely elementary and self-contained. In particular, we use:
\begin{enumerate}[label=(\roman*)]
\item a continuous prime-power ramp kernel $(y/n - 1)_+$ whose Mellin transform is $(s(s-1))^{-1}$, explicitly desingularized at $s=1$ via divided differences;
\item a strictly pointwise-to-$L^1$ one-sided integral inequality on $[U, V] = [\log Y, \lambda \log Y]$ paired with a Stirling-free two-sided Chernoff tail bound on $(0, \infty) \setminus [U, V]$; and
\item a purely geometric vertical equidistance-breaking perturbation that isolates a strictly nearest pole on a local Cauchy disk inside $\{\Re(s) > 0\}$, avoiding both global critical-strip contour shifts and Dirichlet/Bohr simultaneous Diophantine approximation.
\end{enumerate}

\section{Admissible Character Families and Centered Ramp Sums}
\label{sec:02:ramps}

\begin{definition}[Admissible Dirichlet character and centered ramp sum]
\label{def:02:ramp}
A pair $(N, \chi)$ consisting of a modulus $N \ge 1$ and a Dirichlet character $\chi \pmod{N}$ is called \emph{admissible} if either $N = 1$ (so $\chi = \mathbf{1}$ is the trivial character of modulus $1$ and $L(s,\chi) = \zeta(s)$) or $\chi \neq \mathbf{1}$ is a nonprincipal character modulo $N$. For an admissible pair $(N, \chi)$ and any real $y > 0$, we define the \emph{raw ramp sum}
\begin{equation}
\label{eq:02:raw-ramp}
  R(\chi, y) \;:=\; \sum_{1 \le n \le y} \Lambda(n)\chi(n)\!\left(\frac{y}{n} - 1\right)
  \;=\; \sum_{n=1}^{\infty} \Lambda(n)\chi(n)\max\!\left(\frac{y}{n} - 1,\, 0\right)
\end{equation}
and the \emph{centered ramp sum}
\begin{equation}
\label{eq:02:centered-ramp}
  G(\chi, y) \;:=\;
  \begin{cases}
    \displaystyle R(\mathbf{1}, y) - y\log y + (1 + \gamma)y, & N = 1\ (\chi = \mathbf{1}), \\[6pt]
    \displaystyle R(\chi, y) + y\,\frac{L'}{L}(1, \chi), & \chi \neq \mathbf{1},
  \end{cases}
\end{equation}
where $\Lambda$ is the von~Mangoldt function and $\gamma = \lim_{M\to\infty}(\sum_{n=1}^M n^{-1} - \log M)$ is Euler's constant.
\end{definition}
\leanref{src/Formalism/StatementOC.lean:13}{Formalism.Gz}
\leanref{src/Formalism/SharpRestriction/RamificationFrobenius/TotientGaloisClasses/Analytic.lean:23}{Formalism.HC.An.gChi}

\begin{remark}[Necessity of admissibility]
If $\chi = \mathbf{1} \pmod{N}$ is principal with $N > 1$, then $\chi(n) = 0$ for $\gcd(n, N) > 1$, so $R(\chi, y)$ differs from $R(\mathbf{1}, y)$ by $\sum_{p \mid N}\sum_{p^k \le y} (\log p)(y/p^k - 1) = y\sum_{p \mid N}\frac{\log p}{p-1} + O_N(\log y)$, which shifts the linear term at $s=1$. Requiring admissibility guarantees that the principal character enters only at modulus $N=1$, matching the exact principal part of $-\zeta'(s)/\zeta(s)$ at $s=1$.
\end{remark}

\begin{definition}[Positively weighted admissible character family]
\label{def:02:char-family}
A \emph{positively weighted admissible character family} $\mathcal{F} = (I, (N_i)_{i\in I}, (\chi_i)_{i\in I}, (a_i)_{i\in I})$ consists of a finite index set $I$, admissible Dirichlet characters $\chi_i \pmod{N_i}$, and strictly positive real weights $a_i > 0$ ($i \in I$). We associate to $\mathcal{F}$:
\begin{enumerate}[label=\textup{(\roman*)}]
\item the real-valued \emph{family ramp sum} $G_{\mathcal{F}}(y) := \sum_{i \in I} a_i \Re G(\chi_i, y)$ ($y > 0$) and its logarithmic pull-back $g_{\mathcal{F}}(u) := G_{\mathcal{F}}(e^u)$ ($u \in \RR$); and
\item the \emph{right half-plane zero set}
  \begin{equation}
  \label{eq:02:zero-set}
    Z(\mathcal{F}) \;:=\; \bigl\{\rho \in \CC : \Re(\rho) > 0 \text{ and } \exists\, i \in I \text{ with } L(\rho, \chi_i) = 0 \text{ or } L(\bar{\rho}, \chi_i) = 0\bigr\}.
  \end{equation}
\end{enumerate}
\end{definition}
\leanref{src/Development/AutomorphicArithmetic/DedekindLFunctions/CharFamily.lean:24}{Deduction.CharFamily}

Notice that the cyclic family $\mathcal{F}_\psi$ of \cref{ch:01} ($I = \langle \psi \rangle$, $a_{\mathbf{1}} = 1$, $a_\xi = (3o+1)^{-1}$ for $\xi \neq \mathbf{1}$) is a positively weighted admissible family whose ramp sum is $G_{\mathcal{F}_\psi}(y) = G_\psi(y)$.

We first record an elementary quadratic growth bound for $G_{\mathcal{F}}(y)$ that uses neither the prime number theorem nor Chebyshev's estimate.

\begin{lemma}[Elementary quadratic growth and continuity]
\label{lem:02:growth}
For every admissible pair $(N, \chi)$, the function $y \mapsto G(\chi, y)$ is continuous on $[1, \infty)$ and satisfies $|R(\chi, y)| \le 2y^2$ for all $y \ge 1$. Consequently, for every positively weighted admissible family $\mathcal{F}$, the function $g_{\mathcal{F}}(u) = G_{\mathcal{F}}(e^u)$ is continuous on $[0, \infty)$ and there exists an explicit constant $C_0 = C_0(\mathcal{F}) > 0$ such that
\begin{equation}
\label{eq:02:quad-growth}
  |g_{\mathcal{F}}(u)| \;\le\; C_0 e^{2u} \qquad \text{for all } u \ge 0.
\end{equation}
\end{lemma}

\begin{proof}
Each term $\Lambda(n)\chi(n)\max(y/n - 1, 0)$ is continuous in $y$ and vanishes for $n > y$, so $y \mapsto R(\chi, y)$ is a finite sum of continuous functions on any bounded interval $[1, Y]$, hence continuous on $[1,\infty)$. For $1 \le n \le y$, we have $\Lambda(n) \le \log n \le n$ and $0 \le y/n - 1 \le y/n$, whence $\Lambda(n)(y/n - 1) \le y$. Since there are at most $\lfloor y \rfloor \le y$ positive integers $n \le y$, we immediately obtain
\[
  |R(\chi, y)| \;\le\; \sum_{1 \le n \le y} \Lambda(n)\!\left(\frac{y}{n} - 1\right) \;\le\; y \cdot \lfloor y \rfloor \;\le\; y^2 \;\le\; 2y^2 \qquad (y \ge 1).
\]
For $N = 1$, $|y\log y| \le y^2$ and $(1+\gamma)y \le 2y^2$ on $[1, \infty)$; for $\chi \neq \mathbf{1}$, $|y\,\frac{L'}{L}(1,\chi)| \le |\frac{L'}{L}(1,\chi)| y^2$. Summing over $i \in I$ with weights $a_i > 0$ and setting $y = e^u$ ($u \ge 0$) yields \eqref{eq:02:quad-growth} with
\[
  C_0(\mathcal{F}) \;:=\; \sum_{i \in I} a_i \left(2 + \begin{cases} 3, & N_i = 1, \\ \bigl|\frac{L'}{L}(1,\chi_i)\bigr|, & \chi_i \neq \mathbf{1} \end{cases}\right). \qedhere
\]
\end{proof}

\section{Mellin--Laplace Continuation and Non-Cancelling Residues}
\label{sec:02:laplace}

We now compute the one-sided Laplace transform of $G(\chi, e^u)$ for $\Re(s) > 2$, establish its meromorphic continuation to $\CC \setminus \{0\}$ with holomorphicity at $s = 1$, and construct the symmetrized Laplace transform $\Phi(s)$ of $g_{\mathcal{F}}(u)$ with guaranteed non-cancelling residues and strictly positive real-pole residues.

\begin{lemma}[Mellin--Laplace transform of the centered ramp sum]
\label{lem:02:mellin-ramp}
Let $(N, \chi)$ be an admissible Dirichlet character. For every $s \in \CC$ with $\Re(s) > 2$, the integral $\int_0^\infty G(\chi, e^u) e^{-su}\,du$ converges absolutely and satisfies
\begin{equation}
\label{eq:02:mellin-id}
  \int_0^\infty G(\chi, e^u) e^{-su}\,du \;=\;
  \begin{cases}
    \displaystyle -\frac{\frac{\zeta'(s)}{\zeta(s)} + \frac{1}{s-1} + 1 - (1+\gamma)s}{s(s-1)}, & N = 1\ (\chi = \mathbf{1}), \\[12pt]
    \displaystyle -\frac{\frac{L'(s,\chi)}{L(s,\chi)} - s\,\frac{L'(1,\chi)}{L(1,\chi)}}{s(s-1)}, & \chi \neq \mathbf{1}.
  \end{cases}
\end{equation}
\end{lemma}
\leanref{src/Development/AutomorphicArithmetic/DedekindLFunctions/LandauTheorems/MellinRamp.lean:10}{Deduction.mellin\_ramp}

\begin{proof}
Fix $s = \sigma + it$ with $\sigma > 2$. By \cref{lem:02:growth}, $|G(\chi, e^u) e^{-su}| \le C e^{(2-\sigma)u}$, which is integrable on $[0, \infty)$. For every integer $n \ge 1$, substituting $y = e^u$ ($e^{-su}\,du = y^{-s-1}\,dy$) gives
\[
  \int_0^\infty \max\!\left(\frac{e^u}{n} - 1,\, 0\right) e^{-su}\,du
  \;=\; \int_n^\infty \left(\frac{y}{n} - 1\right) y^{-s-1}\,dy
  \;=\; \frac{1}{n}\left[\frac{y^{1-s}}{1-s}\right]_n^\infty - \left[\frac{y^{-s}}{-s}\right]_n^\infty
  \;=\; \frac{n^{-s}}{s(s-1)}.
\]
To justify termwise summation over $n \ge 1$, observe by Tonelli's theorem that
\[
  \sum_{n=1}^{\infty} \int_0^\infty \Lambda(n)\max\!\left(\frac{e^u}{n} - 1,\, 0\right) e^{-\sigma u}\,du
  \;=\; \sum_{n=1}^{\infty} \frac{\Lambda(n) n^{-\sigma}}{\sigma(\sigma - 1)} \;<\; \infty,
\]
since $\Lambda(n) \le \log n$ and $\sigma > 2$. Hence Fubini's theorem yields
\[
  \int_0^\infty R(\chi, e^u) e^{-su}\,du \;=\; \frac{1}{s(s-1)}\sum_{n=1}^\infty \Lambda(n)\chi(n) n^{-s}
  \;=\; -\frac{1}{s(s-1)}\frac{L'(s,\chi)}{L(s,\chi)} \qquad (\Re(s) > 2).
\]
For the centering terms, direct integration on $[0, \infty)$ for $\Re(s) > 1$ gives
\[
  \int_0^\infty e^u \cdot e^{-su}\,du \;=\; \frac{1}{s-1}, \qquad
  \int_0^\infty u e^u \cdot e^{-su}\,du \;=\; \frac{1}{(s-1)^2}.
\]
When $\chi \neq \mathbf{1}$, adding $\frac{L'}{L}(1,\chi)\cdot \frac{1}{s-1} = \frac{s\,\frac{L'}{L}(1,\chi)}{s(s-1)}$ to $-\frac{1}{s(s-1)}\frac{L'}{L}(s,\chi)$ gives the second case of \eqref{eq:02:mellin-id}. When $N = 1$ ($\chi = \mathbf{1}$), subtracting $\frac{1}{(s-1)^2}$ and adding $\frac{1+\gamma}{s-1}$ gives
\[
  -\frac{\zeta'(s)/\zeta(s)}{s(s-1)} - \frac{1}{(s-1)^2} + \frac{1+\gamma}{s-1}
  \;=\; -\frac{\frac{\zeta'(s)}{\zeta(s)} + \frac{s}{s-1} - (1+\gamma)s}{s(s-1)}
  \;=\; -\frac{\frac{\zeta'(s)}{\zeta(s)} + \frac{1}{s-1} + 1 - (1+\gamma)s}{s(s-1)},
\]
since $\frac{s}{s-1} = \frac{1}{s-1} + 1$.
\end{proof}

To handle both $\chi = \mathbf{1}$ and $\chi \neq \mathbf{1}$ uniformly at $s = 1$, we introduce the divided-difference desingularization.

\begin{proposition}[Unified meromorphic continuation and symmetrized pole data]
\label{prop:02:pole-data}
Let $\mathcal{F} = (I, (N_i), (\chi_i), (a_i))$ be a positively weighted admissible character family.
\begin{enumerate}[label=\textup{(\roman*)}]
\item For each admissible pair $(N, \chi)$, let $(g_\chi, \kappa_\chi)$ be $((s-1)\zeta(s), 1)$ if $N=1$, and $(L(s,\chi), 0)$ if $\chi \neq \mathbf{1}$. Then $g_\chi$ is entire with $g_\chi(1) \neq 0$, and the auxiliary meromorphic function
  \begin{equation}
  \label{eq:02:nchi-def}
    N_\chi(s) \;:=\; \frac{g_\chi'(s)}{g_\chi(s)} - s\,\frac{g_\chi'(1)}{g_\chi(1)} + \kappa_\chi(1 - s)
  \end{equation}
  is holomorphic at $s = 1$ with $N_\chi(1) = 0$. Consequently, the single-character Laplace transform
  \begin{equation}
  \label{eq:02:phichi-def}
    \Phi_\chi(s) \;:=\; -\frac{1}{s}\,\frac{N_\chi(s)}{s-1} \quad (s \neq 1), \qquad \Phi_\chi(1) \;:=\; -N_\chi'(1),
  \end{equation}
  agrees with $\int_0^\infty G(\chi, e^u) e^{-su}\,du$ for $\Re(s) > 2$, extends meromorphically to $\CC \setminus \{0\}$, is holomorphic at $s = 1$ and at every $s \in \CC \setminus \{0\}$ with $L(s, \chi) \neq 0$, and at every zero $\rho \in \CC \setminus \{0, 1\}$ of $L(s, \chi)$ of multiplicity $m_\chi(\rho) = \ord_\rho L(\cdot, \chi) \ge 1$ has a simple pole with exact principal part
  \[
    \Phi_\chi(s) \;=\; \frac{m_\chi(\rho)/\bigl(\rho(1-\rho)\bigr)}{s - \rho} + h_\chi(s),
  \]
  where $h_\chi$ is holomorphic near $\rho$.

\item Define the \emph{symmetrized Laplace transform} of $\mathcal{F}$ by
  \begin{equation}
  \label{eq:02:phi-sym}
    \Phi(s) \;:=\; \sum_{i \in I} \frac{a_i}{2}\bigl(\Phi_{\chi_i}(s) + \Phi_{\chi_i^{-1}}(s)\bigr).
  \end{equation}
  Then:
  \begin{enumerate}[label=\textup{(\alph*)}]
  \item For every $s \in \CC$ with $\Re(s) > 2$, $\Phi(s) = \int_0^\infty g_{\mathcal{F}}(u) e^{-su}\,du$.
  \item $\Phi$ is meromorphic on $\CC \setminus \{0\}$, holomorphic on $\{\Re(s) > 0\} \setminus Z(\mathcal{F})$ (including at $s=1$), and every $\rho \in Z(\mathcal{F})$ satisfies $0 < \Re(\rho) < 1$ (in particular $Z(\mathcal{F})$ is discrete in $\{\Re(s) > 0\}$).
  \item At every $\rho \in Z(\mathcal{F})$, $\Phi$ has a simple pole with nonzero residue
    \begin{equation}
    \label{eq:02:residue}
      c(\rho) \;:=\; \operatorname{Res}_{s=\rho}\Phi(s) \;=\; \frac{S(\rho)}{\rho(1-\rho)} \;\neq\; 0,
      \qquad S(\rho) \;:=\; \sum_{i \in I} \frac{a_i}{2}\bigl(m_{\chi_i}(\rho) + m_{\chi_i^{-1}}(\rho)\bigr) \;>\; 0.
    \end{equation}
  \item \emph{(Real-pole positivity.)} If $\beta \in Z(\mathcal{F}) \cap \RR$, then $\beta \in (0, 1)$ and the residue $c(\beta) = \frac{S(\beta)}{\beta(1-\beta)}$ is \emph{strictly positive and real}.
  \end{enumerate}
\end{enumerate}
\end{proposition}
\leanref{src/Development/AutomorphicArithmetic/DedekindLFunctions/LandauTheorems/poleData.lean:248}{Deduction.poleData}

\begin{proof}
(i) For $N = 1$, $g_{\mathbf{1}}(s) = (s-1)\zeta(s)$ is entire with $g_{\mathbf{1}}(1) = 1$ and Laurent expansion $\zeta(s) = (s-1)^{-1} + \gamma + O(s-1)$, whence $g_{\mathbf{1}}'(1)/g_{\mathbf{1}}(1) = \gamma$ and
\[
  \frac{g_{\mathbf{1}}'(s)}{g_{\mathbf{1}}(s)} \;=\; \frac{1}{s-1} + \frac{\zeta'(s)}{\zeta(s)}.
\]
Thus $N_{\mathbf{1}}(s) = \frac{\zeta'(s)}{\zeta(s)} + \frac{1}{s-1} - \gamma s + (1-s) = \frac{\zeta'(s)}{\zeta(s)} + \frac{1}{s-1} + 1 - (1+\gamma)s$, matching the numerator of \eqref{eq:02:mellin-id} identically. For $\chi \neq \mathbf{1}$, $L(s, \chi)$ is entire and $L(1, \chi) \neq 0$ by Dirichlet's nonvanishing theorem (\litref{Iwaniec--Kowalski (2004, Theorem~4.9)~\cite{IwaniecKowalski2004}}), so $g_\chi(1) \neq 0$ and $N_\chi(s) = \frac{L'(s,\chi)}{L(s,\chi)} - s\frac{L'(1,\chi)}{L(1,\chi)}$, again matching \eqref{eq:02:mellin-id} identically.

In both cases, $g_\chi(1) \neq 0$ implies $N_\chi$ is holomorphic near $s = 1$ and $N_\chi(1) = \frac{g_\chi'(1)}{g_\chi(1)} - \frac{g_\chi'(1)}{g_\chi(1)} + 0 = 0$. Therefore the divided difference $s \mapsto N_\chi(s)/(s-1)$ has a removable singularity at $s = 1$ with limit $N_\chi'(1)$, so $\Phi_\chi$ is holomorphic at $s = 1$. Away from $s \in \{0, 1\}$, the poles of $\Phi_\chi$ are precisely the zeros of $g_\chi(s)$ (equivalently, the zeros of $L(s, \chi)$). If $L(\rho, \chi) = 0$ with $\rho \notin \{0, 1\}$ and multiplicity $m_\chi(\rho) \ge 1$, then $g_\chi'(s)/g_\chi(s) = \frac{m_\chi(\rho)}{s-\rho} + O(1)$ near $\rho$, so
\[
  \Phi_\chi(s) \;=\; -\frac{1}{s(s-1)}\left(\frac{m_\chi(\rho)}{s-\rho} + O(1)\right)
  \;=\; \frac{-m_\chi(\rho)/(\rho(\rho-1))}{s-\rho} + h_\chi(s)
  \;=\; \frac{m_\chi(\rho)/(\rho(1-\rho))}{s-\rho} + h_\chi(s).
\]

(ii) By Schwarz reflection for the regularized Hurwitz zeta function $\zeta(\bar{s}, a) = \overline{\zeta(s, a)}$, every Dirichlet $L$-function obeys $L(\bar{s}, \chi^{-1}) = \overline{L(s, \chi)}$ on $\CC \setminus \{1\}$. Consequently, $L(\bar{\rho}, \chi_i) = 0 \iff L(\rho, \chi_i^{-1}) = 0$, and $m_{\chi_i^{-1}}(\rho) = m_{\chi_i}(\bar{\rho})$. Likewise, $\overline{G(\chi_i, y)} = G(\chi_i^{-1}, y)$ for all $y > 0$, so
\[
  \Re G(\chi_i, e^u) \;=\; \frac{1}{2}\bigl(G(\chi_i, e^u) + G(\chi_i^{-1}, e^u)\bigr).
\]
Integrating against $e^{-su}\,du$ for $\Re(s) > 2$ and summing with weights $a_i$ proves (a).

For (b), by the classical Hadamard--de la Vall\'ee Poussin theorem (\litref{Iwaniec--Kowalski (2004, Theorem~5.26)~\cite{IwaniecKowalski2004}}), every Dirichlet $L$-function $L(s, \chi)$ is nonvanishing on $\Re(s) \ge 1$. Hence every $\rho \in Z(\mathcal{F})$ satisfies $0 < \Re(\rho) < 1$ (in particular $\rho \notin \{0, 1\}$). Because $I$ is finite and each $g_{\chi_i}, g_{\chi_i^{-1}}$ is a nonzero entire function, $Z(\mathcal{F})$ is a discrete subset of $\{\Re(s) > 0\}$, and $\Phi$ is holomorphic on $\{\Re(s) > 0\} \setminus Z(\mathcal{F})$.

For (c), if $\rho \in Z(\mathcal{F})$, then there exists $i_0 \in I$ such that either $m_{\chi_{i_0}}(\rho) \ge 1$ or $m_{\chi_{i_0}^{-1}}(\rho) = m_{\chi_{i_0}}(\bar{\rho}) \ge 1$. Since every weight satisfies $a_i > 0$ and every multiplicity satisfies $m_{\chi_i}(\rho), m_{\chi_i^{-1}}(\rho) \ge 0$, the sum $S(\rho)$ in \eqref{eq:02:residue} satisfies $S(\rho) \ge a_{i_0}/2 > 0$. Since $\rho \notin \{0, 1\}$, we obtain $c(\rho) = S(\rho)/(\rho(1-\rho)) \neq 0$.

For (d), if $\beta \in Z(\mathcal{F}) \cap \RR$, then (b) yields $\beta \in (0, 1)$, so $\beta(1 - \beta) > 0$ and $c(\beta) = S(\beta)/(\beta(1-\beta)) \in (0, \infty)$ is strictly positive and real.
\end{proof}

\begin{remark}[Sign cross-check]
As an independent sanity check against sign inversions (Safeguard~\textbf{G3} of \cref{sec:00:manifesto}), note that by the explicit formula for $\psi_0(\chi, u) = \sum_{n \le u}\Lambda(n)\chi(n)$, integrating $G(\chi, y) = -y\int_y^\infty (\psi_0(\chi, u) - \text{main}) u^{-2}\,du$ converts a real zero contribution $-u^\beta/\beta$ of $\psi_0$ into $+y^\beta / (\beta(1-\beta))$, whose one-sided Laplace transform is $+\frac{1}{\beta(1-\beta)}(s-\beta)^{-1}$. Thus the positive sign in \eqref{eq:02:residue} is exact.
\end{remark}

\section{One-Sided Window Core Inequality and Stirling-Free Chernoff Tails}
\label{sec:02:chernoff}

We next translate a one-sided pointwise upper bound on $g_{\mathcal{F}}(u)$ over a single logarithmic window $[U, V] \subset (0, \infty)$ into a uniform upper bound on the $k$-th Taylor coefficients of $\Phi(s)$ along the vertical line $\Re(s) = \sigma > 2$. Notice that the argument goes strictly from a \emph{pointwise} upper bound on $[U, V]$ to an $L^1$ integral bound, never in the reverse direction (Safeguard~\textbf{G1}).

\begin{lemma}[Iterated differentiation under the Laplace integral]
\label{lem:02:iterated-laplace}
Let $g : [0, \infty) \to \RR$ be continuous with $|g(u)| \le C_0 e^{2u}$ for all $u \ge 0$, and let $\Phi(s) = \int_0^\infty g(u) e^{-su}\,du$ for $\Re(s) > 2$. Then $\Phi$ is holomorphic on $\{\Re(s) > 2\}$, and for every integer $k \ge 0$ and every $s \in \CC$ with $\Re(s) > 2$,
\begin{equation}
\label{eq:02:iter-lap}
  \frac{(-1)^k \Phi^{(k)}(s)}{k!} \;=\; \int_0^\infty g(u)\,\frac{u^k e^{-su}}{k!}\,du.
\end{equation}
\end{lemma}

\begin{proof}
Fix $s_0 \in \CC$ with $\Re(s_0) > 2$, and set $r := (\Re(s_0) - 2)/2 > 0$. On the disk $B(s_0, r)$, every $s$ satisfies $\Re(s) \ge 2 + r$, so for any $m \ge 0$ we have $|g(u) u^m e^{-su}| \le C_0 u^m e^{-ru}$, which is integrable on $[0, \infty)$ with integral $C_0 m! / r^{m+1}$. By induction on $k$ and dominated convergence for differentiating under the integral sign on $B(s_0, r)$, $\Phi^{(k)}(s) = \int_0^\infty g(u)(-u)^k e^{-su}\,du$, giving \eqref{eq:02:iter-lap}.
\end{proof}

\begin{proposition}[One-sided window core inequality]
\label{prop:02:window-ineq}
Let $g : [0, \infty) \to \RR$ be continuous with $|g(u)| \le C_0 e^{2u}$ for all $u \ge 0$, and let $\Phi(s) = \int_0^\infty g(u) e^{-su}\,du$ for $\Re(s) > 2$. Fix reals $\sigma > 2$, $\theta < \sigma$, an interval $[U, V] \subset (0, \infty)$, and a continuous nonnegative function $R : [U, V] \to [0, \infty)$ such that
\[
  g(u) \;\le\; R(u) \qquad \text{for all } u \in [U, V].
\]
Then for every integer $k \ge 0$ and every $\gamma \in \RR$,
\begin{equation}
\label{eq:02:window-core}
  \left|\frac{(-1)^k \Phi^{(k)}(\sigma + i\gamma)}{k!}\right|
  + \Re\!\left(\frac{(-1)^k \Phi^{(k)}(\sigma)}{k!}\right)
  \;\le\; \frac{2}{(\sigma - \theta)^k}\int_U^V R(u) e^{-\theta u}\,du
  + 2 C_0 \int_{(0,\infty) \setminus [U,V]} \frac{u^k e^{(2-\sigma)u}}{k!}\,du.
\end{equation}
\end{proposition}
\leanref{src/Development/AutomorphicArithmetic/DedekindLFunctions/LandauTheorems/landauWindowIneq.lean:27}{Deduction.landau\_window\_ineq}

\begin{proof}
Set $w_k(u) := \frac{u^k e^{-\sigma u}}{k!} \ge 0$ for $u \ge 0$. Applying \cref{lem:02:iterated-laplace} at $\sigma + i\gamma$ and at $\sigma$, and using that $g(u)$ and $w_k(u)$ are real with $|e^{-(\sigma + i\gamma)u}| = e^{-\sigma u}$, we obtain
\[
  \left|\frac{(-1)^k \Phi^{(k)}(\sigma + i\gamma)}{k!}\right|
  \;\le\; \int_0^\infty |g(u)|\, w_k(u)\,du,
  \qquad
  \Re\!\left(\frac{(-1)^k \Phi^{(k)}(\sigma)}{k!}\right)
  \;=\; \int_0^\infty g(u)\, w_k(u)\,du.
\]
Adding the two inequalities yields
\[
  \left|\frac{(-1)^k \Phi^{(k)}(\sigma + i\gamma)}{k!}\right|
  + \Re\!\left(\frac{(-1)^k \Phi^{(k)}(\sigma)}{k!}\right)
  \;\le\; \int_0^\infty \bigl(|g(u)| + g(u)\bigr)\, w_k(u)\,du.
\]
Split $(0, \infty) = [U, V] \sqcup ((0, \infty) \setminus [U, V])$:
\begin{itemize}
\item On $[U, V]$, since $g(u) \le R(u)$ and $R(u) \ge 0$, we have the \emph{pointwise} bound $|g(u)| + g(u) = 2\max(g(u), 0) \le 2R(u)$. Moreover, setting $a := \sigma - \theta > 0$, the elementary power-series inequality $(au)^k / k! \le e^{au}$ ($u \ge 0$) gives
  \[
    w_k(u) \;=\; \frac{u^k e^{-(\sigma-\theta)u}}{k!}\, e^{-\theta u} \;\le\; \frac{e^{-\theta u}}{(\sigma - \theta)^k}.
  \]
  Integrating over $[U, V]$ gives $\int_U^V (|g(u)| + g(u)) w_k(u)\,du \le \frac{2}{(\sigma-\theta)^k}\int_U^V R(u) e^{-\theta u}\,du$.
\item On $(0, \infty) \setminus [U, V]$, we have $|g(u)| + g(u) \le 2|g(u)| \le 2C_0 e^{2u}$, which gives the second term of \eqref{eq:02:window-core}. \qedhere
\end{itemize}
\end{proof}

The beauty of \cref{prop:02:window-ineq} is that the exterior tail integral is a purely Gamma-type integral whose peak lies at $u \approx k / (\sigma - 2)$. By picking $\sigma$ and $k$ explicitly in terms of $\delta = \lambda - 1 > 0$ and $U$, we place the peak strictly inside the midpoint of $[U, (1+\delta)U]$ and obtain an exact, Stirling-free two-sided Chernoff bound that decays as $O(\sigma^{-k})$—strictly smaller than $(\sigma - \theta)^{-k}$ whenever $\theta > 0$!

\begin{proposition}[Explicit Stirling-free two-sided Chernoff tail bound]
\label{prop:02:chernoff}
Let $\delta > 0$ (so $\lambda := 1 + \delta > 1$) and define the explicit rational constants
\begin{equation}
\label{eq:02:chernoff-params}
  l \;:=\; 2 + \frac{6}{\delta}, \qquad
  M \;:=\; l(1 + \delta), \qquad
  \sigma \;:=\; M(2 + l), \qquad
  C_{\mathrm{tail}}(\delta) \;:=\; \frac{e}{\sigma - 2 - l} + \frac{1}{\sigma - 2 + l}.
\end{equation}
Then $\sigma > 4 > 2$ and $\sigma - 2 - l > 0$. Moreover, for every real $U \ge 1$, setting $V := (1 + \delta)U$ and $k := \lceil (1 + \delta/2)\sigma U \rceil$, we have $k \ge U$ and
\begin{equation}
\label{eq:02:chernoff-bound}
  \int_{(0, \infty) \setminus [U, V]} \frac{u^k e^{(2-\sigma)u}}{k!}\,du \;\le\; \frac{C_{\mathrm{tail}}(\delta)}{\sigma^k}.
\end{equation}
\end{proposition}
\leanref{src/Development/AutomorphicArithmetic/DedekindLFunctions/LandauTheorems/landauTails.lean:18}{Deduction.landau\_tails\_chernoff}

\begin{proof}
Observe first that $l > 2$, $M = l + l\delta = 2 + \frac{6}{\delta} + 2\delta + 6 > 8$, and $\sigma = M(2+l) > 8 \cdot 4 = 32 > 4$. Also,
\[
  a \;:=\; \sigma - 2 - l \;=\; (M - 1)(2 + l) \;>\; 0, \qquad
  b \;:=\; \sigma - 2 + l \;>\; \sigma \;>\; 0.
\]
Clearly $(1 + \delta/2)\sigma > 1$, so $k = \lceil (1 + \delta/2)\sigma U \rceil \ge U$ for all $U \ge 1$.

\emph{Step~1: Exponential tilting.}
Put $s := \sigma - 2 > l > 0$. On the right tail $u > V = (1+\delta)U$, we have $l(u - V) > 0$, hence $e^{-su} \le e^{-lV} e^{-(s-l)u} = e^{-MU} e^{-au}$ (using $lV = l(1+\delta)U = MU$). On the left tail $0 < u < U$, we have $l(U - u) > 0$, hence $e^{-su} \le e^{lU} e^{-(s+l)u} = e^{lU} e^{-bu}$. Integrating over $(0, \infty)$ using the elementary identity $\int_0^\infty \frac{u^k e^{-\alpha u}}{k!}\,du = \alpha^{-(k+1)}$ ($\alpha > 0$) yields
\begin{equation}
\label{eq:02:tilt-sum}
  \int_{(0,\infty) \setminus [U, V]} \frac{u^k e^{-su}}{k!}\,du
  \;\le\; \frac{e^{-MU}}{a^{k+1}} + \frac{e^{lU}}{b^{k+1}}.
\end{equation}

\emph{Step~2: Elementary change-of-base criterion.}
For any positive real $\alpha$ and reals $x, c$, we claim that
\begin{equation}
\label{eq:02:change-base}
  x + k\,\frac{\sigma - \alpha}{\alpha} \;\le\; c \quad\Longrightarrow\quad \frac{e^x}{\alpha^{k+1}} \;\le\; \frac{e^c}{\alpha}\,\sigma^{-k}.
\end{equation}
Indeed, $\frac{\sigma}{\alpha} = 1 + \frac{\sigma - \alpha}{\alpha} \le \exp\!\left(\frac{\sigma - \alpha}{\alpha}\right)$ by the universal inequality $1 + t \le e^t$ ($t \in \RR$). Raising to the $k$-th power and multiplying by $e^x / (\alpha \sigma^k)$ gives $e^x \alpha^{-(k+1)} = \frac{1}{\alpha\sigma^k} e^x (\sigma/\alpha)^k \le \frac{1}{\alpha\sigma^k} \exp(x + k\frac{\sigma-\alpha}{\alpha}) \le \frac{e^c}{\alpha}\sigma^{-k}$.

\emph{Step~3: Verification of the two exponent inequalities.}
We verify \eqref{eq:02:change-base} for $\alpha = a$ ($c = 1$) and $\alpha = b$ ($c = 0$) using three exact rational identities:
\begin{itemize}
\item First, from $l = 2 + 6/\delta$ we have $\frac{\delta}{2}(l - 2) = 3$, and therefore
  \[
    \left(1 + \frac{\delta}{2}\right)(2 + l) \;=\; 2 + l + \delta + \frac{\delta}{2}l
    \;=\; l(1 + \delta) + 2 + \delta - \frac{\delta}{2}l
    \;=\; M + 2 - \frac{\delta}{2}(l - 2) \;=\; M - 1.
  \]
  Multiplying by $M$ gives $(1 + \delta/2)\sigma = M(M - 1)$.
\item For the right tail ($\alpha = a = (M-1)(2+l)$): we have $\sigma - a = 2 + l$, so $\frac{\sigma - a}{a} = \frac{1}{M - 1}$. Because $k = \lceil (1+\delta/2)\sigma U \rceil < (1+\delta/2)\sigma U + 1 = M(M-1)U + 1$, we obtain
  \[
    -MU + k\,\frac{\sigma - a}{a}
    \;<\; -MU + \frac{M(M-1)U + 1}{M - 1}
    \;=\; \frac{1}{M - 1} \;\le\; 1,
  \]
  since $M > 8 > 2$. Hence \eqref{eq:02:change-base} with $c = 1$ gives $e^{-MU} a^{-(k+1)} \le \frac{e}{a}\sigma^{-k}$.
\item For the left tail ($\alpha = b = \sigma + l - 2$): we have $\sigma - b = -(l - 2) < 0$. Using $\frac{\delta}{2}(l-2) = 3$ again gives $(1 + \delta/2)(l - 2) = l + 1$. Moreover,
  \[
    (l + 1)\sigma - l b \;=\; (l + 1)\sigma - l(\sigma + l - 2) \;=\; \sigma + 2l - l^2.
  \]
  Since $M = l(1+\delta) > l$, we have $\sigma = M(2+l) > l(2+l) = l^2 + 2l$, and thus $\sigma + 2l - l^2 > 4l > 0$, i.e., $(l+1)\sigma / b > l$. Since $k \ge (1+\delta/2)\sigma U$, it follows that
  \[
    l U + k\,\frac{\sigma - b}{b}
    \;=\; l U - k\,\frac{l - 2}{b}
    \;\le\; l U - \frac{(1 + \delta/2)(l - 2)\sigma U}{b}
    \;=\; \left(l - \frac{(l + 1)\sigma}{b}\right) U \;<\; 0.
  \]
  Hence \eqref{eq:02:change-base} with $c = 0$ gives $e^{lU} b^{-(k+1)} \le \frac{1}{b}\sigma^{-k}$.
\end{itemize}
Combining both bounds inside \eqref{eq:02:tilt-sum} yields \eqref{eq:02:chernoff-bound}.
\end{proof}

\begin{example}[Parameters for the active window $\lambda = 3/2$]
For our monograph's locked window exponent $\lambda = 3/2$ ($\delta = 1/2$), \eqref{eq:02:chernoff-params} yields exact integers $l = 14$, $M = 21$, $\sigma = 336$, $a = 320$, $b = 348$, $C_{\mathrm{tail}}(1/2) = \frac{e}{320} + \frac{1}{348} < 0.012$, and $k = \lceil 420 U \rceil$.
\end{example}

Combining \cref{prop:02:window-ineq,prop:02:chernoff} yields the uniform Taylor-coefficient window upper bound along $\Re(s) = \sigma$.

\begin{corollary}[Uniform Taylor-coefficient window bound]
\label{cor:02:window-key}
Let $\mathcal{F}$ be a positively weighted admissible character family, with $g = g_{\mathcal{F}}$, $\Phi$ as in \eqref{eq:02:phi-sym}, and $C_0 = C_0(\mathcal{F})$ from \cref{lem:02:growth}. Fix $\delta > 0$, let $\sigma = \sigma(\delta) > 4$ and $C_{\mathrm{tail}} = C_{\mathrm{tail}}(\delta)$ be given by \eqref{eq:02:chernoff-params}, and fix $\theta \in (0, 1]$ and $A > 0$. Suppose that for arbitrarily large $Y \ge e$,
\begin{equation}
\label{eq:02:window-hyp}
  G_{\mathcal{F}}(y) \;\le\; A\, y^\theta \qquad \text{for all } y \in \bigl[Y,\, Y^{1+\delta}\bigr].
\end{equation}
Set $C' := 2 A + 2 C_0 C_{\mathrm{tail}}$. Then for every $\eta > 0$ and every $U_0 \ge 1$, there exist $U \ge U_0$ and an integer $k \ge U$ such that \emph{uniformly for all $\gamma \in \RR$},
\begin{equation}
\label{eq:02:taylor-key}
  \left|\frac{(-1)^k \Phi^{(k)}(\sigma + i\gamma)}{k!}\right|
  + \Re\!\left(\frac{(-1)^k \Phi^{(k)}(\sigma)}{k!}\right)
  \;\le\; \frac{C' e^{\eta U}}{(\sigma - \theta)^k}.
\end{equation}
\end{corollary}
\leanref{src/Development/AutomorphicArithmetic/DedekindLFunctions/LandauTheorems/landauWindow.lean:40}{Deduction.landau\_window\_key}

\begin{proof}
Choose $U_1 \ge \max(U_0, 1)$ large enough that $\delta U_1 \le e^{\eta U_1}$ (for instance, $U_1 \ge \max(U_0, 1, \delta / \eta)$ suffices since $e^{\eta U} \ge \eta U \ge \delta$ and wait—to get $\delta U \le e^{\eta U}$, note $e^{\eta U} \ge \frac{1}{2}(\eta U)^2 \ge \delta U$ as soon as $U \ge 2\delta/\eta^2$). By hypothesis, pick $Y \ge e^{U_1}$ satisfying \eqref{eq:02:window-hyp}, and set $U := \log Y \ge U_1$, $V := (1+\delta)U$, and $k := \lceil (1+\delta/2)\sigma U \rceil \ge U$.

On $[U, V]$, \eqref{eq:02:window-hyp} reads $g(u) \le A e^{\theta u} =: R(u)$, so $\int_U^V R(u) e^{-\theta u}\,du = A(V - U) = A \delta U \le A e^{\eta U}$. Meanwhile, since $0 < \theta < \sigma$, we have $\sigma^{-k} < (\sigma - \theta)^{-k} \le e^{\eta U} (\sigma - \theta)^{-k}$. Substituting both bounds and \cref{prop:02:chernoff} into \cref{prop:02:window-ineq} yields
\[
  \left|\frac{(-1)^k \Phi^{(k)}(\sigma + i\gamma)}{k!}\right|
  + \Re\!\left(\frac{(-1)^k \Phi^{(k)}(\sigma)}{k!}\right)
  \;\le\; \frac{2A e^{\eta U}}{(\sigma - \theta)^k} + \frac{2C_0 C_{\mathrm{tail}}}{\sigma^k}
  \;\le\; \frac{(2A + 2C_0 C_{\mathrm{tail}}) e^{\eta U}}{(\sigma - \theta)^k}. \qedhere
\]
\end{proof}

\section{Vertical Equidistance-Breaking and the Window Landau Theorem}
\label{sec:02:landau-proof}

In classical proofs of Landau-type theorems for oscillating Dirichlet series, detecting an off-axis pole $\rho_0 \in Z(\mathcal{F})$ with $\Re(\rho_0) > \theta$ typically invokes Dirichlet's pigeonhole theorem or Bohr's almost-periodicity theorem to align phases of multiple poles lying at equal distance from $\sigma + i\gamma$. Here we eliminate simultaneous Diophantine approximation completely: by shifting $\gamma$ slightly off a finite set of perpendicular bisectors, we guarantee a \emph{strictly unique nearest pole} on a local disk inside $\{\Re(s) > 0\}$.

\begin{proposition}[Geometric dichotomy and vertical equidistance-breaking]
\label{prop:02:dichotomy}
Let $Z \subset \{s \in \CC : 0 < \Re(s) < 1\}$ be a discrete subset of $\{\Re(s) > 0\}$, fix $\sigma > 2$ and $\theta \in (0, 1]$, and suppose there exists $\rho_0 \in Z$ with $\Re(\rho_0) > \theta$. Then at least one of the following two cases holds:
\begin{enumerate}[label=\textup{\textbf{Case~\Roman*.}}]
\item \emph{(Off-axis unique nearest pole.)} There exist real numbers $D \in (0, \sigma - \theta]$ and $\gamma \in \RR$, a point $\rho \in Z$, and radii $0 < d < r \le D$ such that, setting $z := \sigma + i\gamma$,
  \[
    Z \cap B(\sigma, D) \;=\; \varnothing, \qquad Z \cap B(z, r) \;=\; \{\rho\}, \qquad |z - \rho| \;=\; d \;<\; r \;\le\; \sigma - \theta.
  \]
\item \emph{(Strictly unique nearest real pole.)} There exist a real pole $\beta \in Z \cap (\theta, 1)$ and radii $0 < d < r \le \sigma - \theta$ such that
  \[
    Z \cap B(\sigma, r) \;=\; \{\beta\}, \qquad \sigma - \beta \;=\; d \;<\; r \;\le\; \sigma - \theta.
  \]
\end{enumerate}
\end{proposition}

\begin{proof}
Consider the closed disk $\overline{B}(\sigma, \sigma - \theta) \subset \{\Re(s) > 0\}$ (which stays in $\{\Re(s) > 0\}$ because $\theta > 0$, so every $s \in \overline{B}(\sigma, \sigma - \theta)$ has $\Re(s) \ge \theta > 0$). Since $Z$ is discrete in $\{\Re(s) > 0\}$ and $\overline{B}(\sigma, \sigma - \theta)$ is compact, the intersection $F_0 := Z \cap \overline{B}(\sigma, \sigma - \theta)$ is finite.

We first establish a clean dichotomy for $F_0$:
\begin{itemize}
\item \emph{Subcase~1:} $F_0 = \varnothing$, or $F_0 \neq \varnothing$ and setting $d_0 := \min_{w \in F_0} |w - \sigma| \le \sigma - \theta$, there exists $\rho_1 \in Z$ with $\sigma - \Re(\rho_1) < d_0$. Put $D := \sigma - \theta$ if $F_0 = \varnothing$ (taking $\rho_1 := \rho_0$, which satisfies $\sigma - \Re(\rho_0) < \sigma - \theta = D$), and put $D := d_0$ otherwise. In both instances we have $0 < D \le \sigma - \theta$, $Z \cap B(\sigma, D) = \varnothing$, and $\sigma - \Re(\rho_1) < D$.
\item \emph{Subcase~2:} $F_0 \neq \varnothing$, $d_0 := \min_{w \in F_0} |w - \sigma| \le \sigma - \theta$, and \emph{every} $w \in Z$ satisfies $d_0 \le \sigma - \Re(w)$. Taking $w = \rho_0$ gives $d_0 \le \sigma - \Re(\rho_0) < \sigma - \theta$. Let $\beta \in F_0$ be any minimizer of $|w - \sigma|$ on $F_0$. Then
  \[
    (\sigma - \Re(\beta))^2 + (\Im(\beta))^2 \;=\; d_0^2 \;\le\; (\sigma - \Re(\beta))^2,
  \]
  forcing $\Im(\beta) = 0$ and $\sigma - \beta = d_0$, hence $\beta = \sigma - d_0 \in (\theta, 1)$ is \emph{strictly unique}! Setting $d := d_0$ and
  \[
    r \;:=\; \min\!\left(\sigma - \theta,\; \min_{w \in F_0 \setminus \{\beta\}} |w - \sigma|\right) \;\in\; (d, \sigma - \theta]
  \]
  (with the inner minimum understood as $+\infty$ if $F_0 = \{\beta\}$) gives $Z \cap B(\sigma, r) = \{\beta\}$, proving \textbf{Case~II}.
\end{itemize}

It remains to deduce \textbf{Case~I} from Subcase~1 via an eight-line vertical equidistance-breaking perturbation. Set $c_0 := \sigma + i\Im(\rho_1)$,
\[
  \varepsilon_0 \;:=\; \frac{1}{2}\min\!\bigl(D - (\sigma - \Re(\rho_1)),\; \sigma - D\bigr) \;>\; 0,
\]
and $F := Z \cap \overline{B}(c_0, D + \varepsilon_0)$. Because $D + \varepsilon_0 < \sigma$, the closed disk $\overline{B}(c_0, D + \varepsilon_0)$ lies strictly inside $\{\Re(s) > 0\}$, so $F$ is finite. For any two distinct points $a, b \in F$ (which satisfy $\Re(a), \Re(b) < 1 < \sigma$), the equidistance equation $|\sigma + i\gamma - a|^2 = |\sigma + i\gamma - b|^2$ expands linearly in $\gamma$ as
\begin{equation}
\label{eq:02:bisector}
  2\gamma\bigl(\Im(b) - \Im(a)\bigr) \;=\; (\sigma - \Re(b))^2 + (\Im(b))^2 - (\sigma - \Re(a))^2 - (\Im(a))^2.
\end{equation}
If $\Im(a) = \Im(b)$, then \eqref{eq:02:bisector} forces $(\sigma - \Re(b))^2 = (\sigma - \Re(a))^2$; since both $\sigma - \Re(a)$ and $\sigma - \Re(b)$ are strictly positive, this implies $\Re(a) = \Re(b)$ and hence $a = b$, contradiction. If $\Im(a) \neq \Im(b)$, then \eqref{eq:02:bisector} has a unique real solution $\gamma_{a,b}$. Therefore the set of bad vertical shifts
\[
  B_{\mathrm{eq}} \;:=\; \bigl\{\gamma_{a,b} : a, b \in F,\; \Im(a) \neq \Im(b)\bigr\}
\]
is \emph{finite}. Choose any $\gamma \in (\Im(\rho_1) - \varepsilon_0, \Im(\rho_1) + \varepsilon_0) \setminus B_{\mathrm{eq}}$ and set $z := \sigma + i\gamma$. Then $|z - c_0| < \varepsilon_0$, so:
\begin{enumerate}[label=(\alph*)]
\item $|z - \rho_1| \le |z - c_0| + |c_0 - \rho_1| < \varepsilon_0 + (\sigma - \Re(\rho_1)) < D$, whence $Z \cap B(z, D)$ is nonempty;
\item by the triangle inequality, $B(z, D) \subset B(c_0, D + \varepsilon_0)$, so $Z \cap B(z, D) \subseteq F$; and
\item because $\gamma \notin B_{\mathrm{eq}}$, the distances $\{|z - w| : w \in F\}$ are pairwise distinct!
\end{enumerate}
Hence the minimizer $\rho \in Z \cap B(z, D)$ of $|z - w|$ is strictly unique. Setting $d := |z - \rho| < D$ and $r := \min\bigl(D,\, \min_{w \in F \setminus \{\rho\}} |z - w|\bigr) \in (d, D]$ yields $Z \cap B(z, r) = \{\rho\}$ and $0 < d < r \le D \le \sigma - \theta$, proving \textbf{Case~I}.
\end{proof}

We now assemble the complete proof of the Window Landau Theorem.

\begin{theorem}[Window Landau Theorem]
\label{thm:02:window-landau}
Let $\mathcal{F} = (I, (N_i), (\chi_i), (a_i))$ be a positively weighted admissible character family (\cref{def:02:char-family}), fix a window exponent $\lambda = 1 + \delta > 1$ ($\delta > 0$), a threshold $\theta \in (0, 1]$, and a constant $A > 0$. Suppose that for every $Y_{\min} \ge 1$ there exists $Y \ge Y_{\min}$ such that
\begin{equation}
\label{eq:02:wl-hyp}
  G_{\mathcal{F}}(y) \;\le\; A\, y^\theta \qquad \text{for all } y \in \bigl[Y,\, Y^{\lambda}\bigr].
\end{equation}
Then for every $i \in I$ and every $s \in \CC$ with $\Re(s) > \theta$, we have $L(s, \chi_i) \neq 0$.
\end{theorem}
\leanref{src/Formalism/SharpRestriction/RamificationFrobenius/TotientGaloisClasses/Window.lean:7}{Formalism.HC.HasWindows}
\leanref{src/Formalism/SharpRestriction/RamificationFrobenius/TotientGaloisClasses/Window.lean:33}{Formalism.HC.Level.FormB}
\leanref{src/Development/AutomorphicArithmetic/DedekindLFunctions/LandauTheorems/ZeroFreeOfRpowWindow.lean:11}{Deduction.zeroFree\_of\_rpow\_window}

\begin{proof}
Suppose for contradiction that there exist $i_0 \in I$ and $\rho_0 \in \CC$ with $\Re(\rho_0) > \theta$ and $L(\rho_0, \chi_{i_0}) = 0$. Then $\rho_0 \in Z(\mathcal{F})$. Let $\Phi$ be the symmetrized Laplace transform \eqref{eq:02:phi-sym} of $\mathcal{F}$, and let $\sigma = \sigma(\delta) > 4$ and $C' = 2A + 2C_0(\mathcal{F})C_{\mathrm{tail}}(\delta) > 0$ be given by \cref{prop:02:chernoff,cor:02:window-key}. By \cref{prop:02:pole-data}(ii)(b), $Z(\mathcal{F}) \subset \{0 < \Re(s) < 1\}$ is discrete in $\{\Re(s) > 0\}$, so \cref{prop:02:dichotomy} applies to $Z := Z(\mathcal{F})$.

Set $e := (d + r)/2 \in (d, r)$ in whichever of Case~I or Case~II occurs, so that
\[
  0 \;<\; d \;<\; e \;<\; r \;\le\; \sigma - \theta.
\]
We derive a unified lower bound on the left-hand side of \eqref{eq:02:taylor-key} via local Cauchy estimates on closed disks strictly inside $\{\Re(s) > 0\}$:

\begin{itemize}
\item \emph{In \textbf{Case~I}:} We have $Z \cap B(z, r) = \{\rho\}$ with $|z - \rho| = d$ and $Z \cap B(\sigma, D) = \varnothing$ with $e < r \le D$. By \cref{prop:02:pole-data}(ii)(c), the desingularized function $H(s) := \Phi(s) - \frac{c(\rho)}{s-\rho}$ is holomorphic on $B(z, r)$, where $c(\rho) \neq 0$, while $\Phi$ itself is holomorphic on $B(\sigma, D)$. By Cauchy's integral formula on the circles $|s - z| = e < r$ and $|s - \sigma| = e < D$,
  \[
    \left|\frac{(-1)^k H^{(k)}(z)}{k!}\right| \;\le\; \frac{M_1}{e^k}, \qquad
    \left|\frac{(-1)^k \Phi^{(k)}(\sigma)}{k!}\right| \;\le\; \frac{M_2}{e^k},
  \]
  where $M_1 := \max_{|s-z|=e} |H(s)|$ and $M_2 := \max_{|s-\sigma|=e} |\Phi(s)|$ are finite constants independent of $k$. Since $\frac{(-1)^k}{k!}\frac{d^k}{ds^k}\bigl(\frac{c(\rho)}{s-\rho}\bigr)\big|_{s=z} = \frac{c(\rho)}{(z-\rho)^{k+1}}$, the reverse triangle inequality gives
  \[
    \left|\frac{(-1)^k \Phi^{(k)}(z)}{k!}\right| \;\ge\; \frac{|c(\rho)|/d}{d^k} - \frac{M_1}{e^k}, \qquad
    \Re\!\left(\frac{(-1)^k \Phi^{(k)}(\sigma)}{k!}\right) \;\ge\; -\frac{M_2}{e^k}.
  \]
  Adding these two estimates and setting $c_{\mathrm{pos}} := |c(\rho)|/d > 0$, $M_{\mathrm{tot}} := M_1 + M_2 \ge 0$, we obtain for every integer $k \ge 0$:
  \begin{equation}
  \label{eq:02:lower-taylor}
    \frac{c_{\mathrm{pos}}}{d^k} \;\le\; \frac{M_{\mathrm{tot}}}{e^k}
    + \left|\frac{(-1)^k \Phi^{(k)}(\sigma + i\gamma)}{k!}\right|
    + \Re\!\left(\frac{(-1)^k \Phi^{(k)}(\sigma)}{k!}\right).
  \end{equation}

\item \emph{In \textbf{Case~II}:} We have $Z \cap B(\sigma, r) = \{\beta\}$ with $\beta \in (\theta, 1)$ real and $\sigma - \beta = d < e < r$. By \cref{prop:02:pole-data}(ii)(d), the residue $c(\beta)$ is \emph{strictly positive and real}. Setting $H(s) := \Phi(s) - \frac{c(\beta)}{s-\beta}$ (holomorphic on $B(\sigma, r)$) and $M_1 := \max_{|s-\sigma|=e} |H(s)|$, we have $(\sigma - \beta)^{k+1} = d^{k+1} > 0$, hence
  \[
    \Re\!\left(\frac{(-1)^k \Phi^{(k)}(\sigma)}{k!}\right)
    \;=\; \frac{c(\beta)/d}{d^k} + \Re\!\left(\frac{(-1)^k H^{(k)}(\sigma)}{k!}\right)
    \;\ge\; \frac{c(\beta)/d}{d^k} - \frac{M_1}{e^k}.
  \]
  Since $|\frac{(-1)^k \Phi^{(k)}(\sigma)}{k!}| \ge 0$, taking $\gamma := 0$, $c_{\mathrm{pos}} := c(\beta)/d > 0$, and $M_{\mathrm{tot}} := M_1$ gives the exact same inequality \eqref{eq:02:lower-taylor}!
\end{itemize}

To derive the final contradiction, put $L := \log(e / d) > 0$ and $\eta := L / 2 > 0$, and choose
\[
  U_0 \;:=\; \max\!\left(1,\; \frac{2(M_{\mathrm{tot}} + C')}{c_{\mathrm{pos}} L}\right).
\]
By \cref{cor:02:window-key}, there exist $U \ge U_0$ and an integer $k \ge U$ such that \eqref{eq:02:taylor-key} holds at $\gamma$. Combining \eqref{eq:02:lower-taylor} with \eqref{eq:02:taylor-key} and using $e < \sigma - \theta$ (so $(\sigma - \theta)^{-k} < e^{-k}$) yields
\[
  \frac{c_{\mathrm{pos}}}{d^k} \;\le\; \frac{M_{\mathrm{tot}}}{e^k} + \frac{C' e^{\eta U}}{(\sigma - \theta)^k}
  \;\le\; \frac{M_{\mathrm{tot}} + C'}{e^k}\, e^{\eta U}.
\]
Multiplying both sides by $e^k e^{-\eta U}$ and noting that $(e/d)^k = e^{k L} \ge e^{U L}$ (because $k \ge U$ and $L > 0$), we find
\[
  c_{\mathrm{pos}}\, e^{(L/2)U} \;=\; c_{\mathrm{pos}}\, e^{(L - \eta)U} \;\le\; c_{\mathrm{pos}}\, e^{kL - \eta U} \;\le\; M_{\mathrm{tot}} + C'.
\]
However, by the elementary bound $e^t \ge 1 + t > t$ ($t > 0$) and our choice of $U \ge U_0$,
\[
  c_{\mathrm{pos}}\, e^{(L/2)U} \;>\; c_{\mathrm{pos}} \cdot \frac{L}{2} U \;\ge\; c_{\mathrm{pos}} \cdot \frac{L}{2} \cdot \frac{2(M_{\mathrm{tot}} + C')}{c_{\mathrm{pos}} L} \;=\; M_{\mathrm{tot}} + C',
\]
a contradiction. Thus $Z(\mathcal{F})$ contains no point with $\Re(\rho) > \theta$, which proves $L(s, \chi_i) \neq 0$ for all $i \in I$ and all $\Re(s) > \theta$.
\end{proof}

\begin{remark}[Complete elimination of Diophantine approximation and critical-strip contours]
Observe that the proof of \cref{thm:02:window-landau} uses no Stirling approximation, no global contour shift across the critical strip, no zero-counting formula $N(T, \chi)$, and no Dirichlet/Bohr alignment of arguments. Every single step rests solely on Cauchy's integral formula on a local disk inside $\{\Re(s) > 0\}$, the elementary inequality $1 + t \le e^t$, and the finiteness of the set of perpendicular bisectors of a finite set of poles.
\end{remark}
